\subsection{\texorpdfstring{\(L^{1}(G)\) 的闭理想}{L\^{}\{1\}(G) 的闭理想}}

Banach 代数 \(L^{1}(G)\) 上的 Fourier 余变换与 \(L^{1}(G)\) 的 Gelfand 变换等同（II，第 209 页）。在这一等同下，回顾：若 I 是 \(L^{1}(G)\) 的一个理想，我们用 V(I) 表示 \(\widehat{G}\) 中由这样的特征标 \(\chi \in \widehat{G}\) 所组成的闭集：对每个函数 \(f \in I\)，f 的 Fourier 余变换在 \(\chi\) 处为零（参看 I，第 30 页）。对 \(\widehat{G}\) 的每个子集 M，我们用 \(\Upsilon(M)\) 表示由 \(f \in L^{1}(G)\)（\(\overline{\mathcal{F}}_{\mathrm{G}}(f)\) 在 M 上为零）所组成的闭理想（I，第 30 页）。

由 II，第 219 页命题 2，Banach 代数 L \(^{1}\) (G) 是正则的。因此，由 I，第 88 页的 §5，我们导出 Fourier 变换与余变换的下列性质：

\begin{enumerate}
\def\labelenumi{\arabic{enumi})}
\item
  若 F 是 \(\widehat{\mathbf{G}}\) 的一个闭子集，K 是 \(\widehat{\mathbf{G}}\) 的一个紧子集，且 \(\mathrm{F} \cap \mathrm{K} = \varnothing\)，则存在函数 \(f \in \mathrm{L}^{1}(\mathrm{G})\)，使 \(\mathcal{F}_{\mathrm{G}}(f)\) 在 F 上等于 0、在 K 上等于 1（I，第 88 页，命题 1；关于此事实及以下各条，可借助 II，第 207 页公式 (8) 从 Fourier 余变换过渡到 Fourier 变换）。
\item
  设 M 是 \(\widehat{G}\) 的一个闭子集。\(L^{1}(G)\) 中满足 \(V(I) = M\) 的理想 I 所成的集合，以 \(\Upsilon(M)\) 为其最大元，而其最小元则是由 Fourier 余变换具有与 M 不相交的紧支集的那些 \(f \in \mathrm{L}^{1}(\mathrm{G})\) 所组成的集合（I，第 91 页，命题 4）。
\item
  设 I 是 \(L^{1}(G)\) 的一个理想，\(g\colon\widehat{G}\to C\) 是一个连续函数。假设对每个 \(\chi\in\widehat{G}\)，存在函数 \(f_{\chi}\in I\)，使 g 与 \(\mathcal{F}_{\mathrm{G}}(f_{\chi})\) 在 \(\chi\) 的一个邻域内相等。此外再假设存在函数 \(f_{\infty}\in I\)，使 g 与 \(\mathcal{F}_{\mathrm{G}}(f_{\infty})\) 在 \(\widehat{G}\) 的一个紧子集之外相等，而当 G 离散时后一条件总满足。则存在函数 \(f\in I\) 使 \(g=\mathcal{F}_{\mathrm{G}}(f)\)（I，第 91 页，推论 2）。
\end{enumerate}

引理 1.　\(\mathrm{L}^{1}(\mathrm{G})\) 中 Fourier 变换具有紧支集的函数所成的空间在 \(\mathrm{L}^{1}(\mathrm{G})\) 中稠密。

由于 \(\mathcal{K}(\widehat{\mathrm{G}})\) 在 \(\mathrm{L}^2 (\widehat{\mathrm{G}})\) 中稠密，且 \(\mathrm{L}^2 (\mathrm{G})\) 的 Fourier 变换是 \(\mathrm{L}^2 (\widehat{\mathrm{G}})\) 上的等距（II，第 215 页，定理 1），\(\mathrm{L}^2 (\mathrm{G})\) 中由 \(f\in \mathrm{L}^2 (\mathrm{G})\)（使 \(\mathcal{F}_{\mathrm{G}}(f)\in \mathcal{K}(\widehat{\mathrm{G}})\)）所组成的子空间 V 在 \(\mathrm{L}^2 (\mathrm{G})\) 中稠密。

设 \(g \in \mathrm{L}^{1}(\mathrm{G})\)。存在 \(g_{1}, g_{2} \in \mathrm{L}^{2}(\mathrm{G})\) 使 \(g = g_{1}g_{2}\)（例如可取 \(g_{1} = |g|^{1/2}\)，并取 \(g_{2}(x) = 0\)（当 \(g(x) = 0\) 时），否则取 \(g_{2}(x) = g(x)/g_{1}(x)\)）。据此我们推出：g 是形如 \(h_{1}h_{2}\) 的函数组成的一个序列的极限，其中 \(h_{1}\) 与 \(h_{2}\) 属于 V。然而 \(\mathcal{F}_{\mathrm{G}}(h_{1}h_{2}) = \mathcal{F}_{\mathrm{G}}(h_{1}) * \mathcal{F}_{\mathrm{G}}(h_{2})\)（II，第 223 页，命题 14），且 \(\mathcal{F}_{\mathrm{G}}(h_{1}) * \mathcal{F}_{\mathrm{G}}(h_{2})\) 属于 \(\mathcal{K}(\mathrm{G})\)。引理得证。

命题 2.　设 I 是 L \(^{1}\) (G) 的一个闭理想，且设 f ∈ L \(^{1}\) (G)。若 \(\overline{\mathcal{F}}_{\text{G}}(f)\) 在 V(I) 的一个邻域上为零，则 f 属于 I。

设 \(\varepsilon > 0\)。存在函数 \(g \in \mathrm{L}^1(\mathrm{G})\) 使 \(\|f - f * g\|_1 < \varepsilon\)（II，第 211 页，命题 8 (iv)）。设 \(h \in \mathrm{L}^1(\mathrm{G})\) 使 \(\overline{\mathcal{F}}_{\mathrm{G}}(h)\) 的支集为紧且 \(\|f\|_1 \|g - h\|_1 < \varepsilon\)（引理 1）。我们有

\[
\| f - f * h \| _ {1} \leqslant \| f - f * g \| _ {1} + \| f * (g - h) \| _ {1} <   2 \varepsilon .
\]

由对 \(f\) 的假设，函数 \(\overline{\mathcal{F}}_{\mathrm{G}}(f*h)=\overline{\mathcal{F}}_{\mathrm{G}}(f)\overline{\mathcal{F}}_{\mathrm{G}}(h)\) 具有紧支集且与 V(I) 不相交，由此推出 \(f*h\in\mathrm{I}\)（上面的注 2）。由于 \(\varepsilon\) 可任意小，我们有 \(f\in\overline{\mathrm{I}}=\mathrm{I}\)。

定理 1.　设 I 是 L \(^{1}\) (G) 的一个异于 L \(^{1}\) (G) 的闭理想。则存在特征标 \(\widehat{x} \in \widehat{\mathrm{G}}\)，使得 \(\mathcal{F}_{\mathrm{G}}(f)(\widehat{x}) = 0\) 对每个 \(f \in \mathrm{I}\) 成立。

由于代数 L \(^{1}\) (G) 既是正则的（II，第 219 页，命题 2）又是半单的（I，第 126 页，命题 22 的推论），且 Fourier 余变换具有紧支集的函数所成的集合在

\(L^{1}(G)\) 中稠密（引理 1），I，第 92 页的推论 1 表明理想 I 含于 \(L^{1}(G)\) 的一个正则极大理想中，即含于某个特征标 \(\widehat{y}\)（\(L^{1}(G)\) 的）的核中（I，第 30 页，定理 2）；于是可取 \(\widehat{x} = \widehat{y}^{-1}\)（参看 II，第 207 页，公式 (8)）。

推论 1（Wiener–Tauber 型定理）

设 \(f \in \mathrm{L}^1(\mathrm{G})\)。若 \(f\) 的 Fourier 变换处处不为零，则函数 \(f * \varepsilon_x \colon g \mapsto f(gx^{-1})\)（其中 \(x\) 取遍 \(\mathrm{G}\)）在 \(\mathrm{L}^1(\mathrm{G})\) 中构成一个完全集（EVT，I，第 12 页，定义 1）。

设 V 是 \(L^{1}(G)\) 的由诸 \(f*\varepsilon_{x}\) 生成的闭向量子空间。由 INT，VIII，§4，命题 20 的推论，空间 V 是 \(L^{1}(G)\) 的一个闭理想。由定理 1 有 \(V = L^{1}(G)\)。

定义 1.　设 g 是 G 上的复值函数，Φ 是 G 上的一个滤子。我们称 g 沿 Φ 是缓变的，如果对每个 ε \textgreater{} 0，存在一个集合 M ∈ Φ 和 G 中 e 的一个邻域 V，使得

\[
 x \in \mathrm{M} \quad \text {and} \quad y \in \mathrm{V} \quad \Longrightarrow \quad | g (x y) - g (x) | \leqslant \varepsilon .
\]

推论 2.　设 \(\Phi\) 是 G 上一个平移不变的滤子。设 \(f \in \mathrm{L}^{1}(\mathrm{G})\) 满足：\(f\) 的 Fourier 变换处处不为零，且 \(\int_{\mathrm{G}} f(x) dx = 1\)。设 \(g \in \mathrm{L}^{\infty}(\mathrm{G})\)。假设 \(f * g\) 有有限极限 \(\alpha\)（沿 \(\Phi\)）。

\begin{enumerate}
\def\labelenumi{\alph{enumi})}
\item
  对每个函数 \(h \in \mathrm{L}^{1}(\mathrm{G})\)（满足 \(\int_{\mathrm{G}} h(x) dx = 1\)），\(h * g\) 沿 \(\Phi\) 的极限等于 \(\alpha\)；
\item
  此外假设 \(g\) 沿 \(\Phi\) 缓变。则 \(g\) 趋于 \(\alpha\)（沿 \(\Phi\)）。
\end{enumerate}

把 \(g\) 换成 \(g - \alpha\)，可化归为 \(\alpha = 0\) 的情形。设 I 是由这样的函数 \(h \in \mathrm{L}^1(\mathrm{G})\) 组成的集合：\(h*g\) 沿 \(\Phi\) 趋于 0。集合 I 是 \(\mathrm{L}^1(\mathrm{G})\) 的一个平移不变向量子空间。它是一个闭向量子空间。事实上，设 \(h \in \bar{\mathrm{I}}\)。对每个函数 \(h_0 \in \mathrm{L}^1(\mathrm{G})\) 及每个 \(x \in \mathrm{G}\)，有

\[
\begin{array}{c} | (h * g) (x) | \leqslant | ((h - h _ {0}) * g) (x) | + | (h _ {0} * g) (x) | \\ \leqslant \| h - h _ {0} \| _ {1} \| g \| _ {\infty} + | (h _ {0} * g) (x) |. \end{array}
\]

对每个 \(\varepsilon > 0\)，存在 \(h_0 \in \mathrm{I}\) 使 \(\|h - h_0\|_1 \|g\|_\infty < \varepsilon\)。取 \(\mathrm{M} \in \Phi\) 使得 \(|(h_0 * g)(x)| < \varepsilon\) 对所有 \(x \in \mathrm{M}\) 成立。于是有 \(|(h * g)(x)| < 2\varepsilon\)（对所有 \(x \in \mathrm{M}\)），因而 \(h * g\) 沿 \(\Phi\) 收敛到 0。这表明 \(h \in \mathrm{I}\)。

因此空间 I 是 \(\mathrm{L}^{1}(\mathrm{G})\) 的一个闭理想。由假设 \(f \in I\)，又因 f 的 Fourier 变换处处不为零，由定理 1 得 \(\mathrm{I} = \mathrm{L}^{1}(\mathrm{G})\)。这就蕴含 a)。

现在假设 \(b\) ) 的条件。设 \(\varepsilon > 0\)。由于 \(g\) 沿 \(\Phi\) 缓变，存在 \(\mathrm{M} \in \Phi\) 及一个紧邻域 \(\mathrm{V}\)（\(e\) 的），使得

\[
 x \in \mathrm{M} \quad \text { and } \quad y \in \mathrm{V} \Longrightarrow | g (y ^ {- 1} x) - g (x) | \leqslant \varepsilon .
\]

设 \(\varphi\) 是 V 的特征函数，\(\mu = \int \varphi(x)dx\)。对每个 \(x \in G\)，有

\[
\frac {1}{\mu} (\varphi * g) (x) = \frac {1}{\mu} \int_ {\mathrm{V}} g (y ^ {- 1} x) d y = g (x) + \frac {1}{\mu} \int_ {\mathrm{V}} (g (y ^ {- 1} x) - g (x)) d y.
\]

因此对所有 \(x \in \mathrm{M}\)，有

\[
\left| \frac {1}{\mu} (\varphi * g) (x) - g (x) \right| \leqslant \varepsilon .
\]

既然由 \(a\) ) 知 \(\varphi * g\) 沿 \(\Phi\) 的极限为零，我们有 \(\limsup_{\Phi} |g| \leqslant \varepsilon\)。由于 \(\varepsilon\) 任意，我们断定 \(g\) 沿 \(\Phi\) 的极限为零，这就证明了 \(b\) )。

引理 2.　设 K 是 G 的一个紧子集。对每个 \(\eta > 0\)，存在函数 \(j \in \mathrm{L}^{1}(\mathrm{G})\)，使得：

\begin{enumerate}
\def\labelenumi{\alph{enumi})}
\item
  \(\|j\|_{1}\leqslant\sqrt{2};\)
\item
  函数 \(\mathcal{F}_{\mathrm{G}}(j)\) 在 \(\widehat{\mathrm{G}}\) 的单位元的一个邻域内等于 1；
\item
  对所有 \(x \in \mathbf{K}\)，有 \(\| j - j * \varepsilon_x\|_1 \leqslant \eta\)。
\end{enumerate}

集合 U \(_{1}\) 由这样的元素 \(\widehat{x} \in \widehat{\mathrm{G}}\) 组成（它们使

\[
| \langle \widehat {x}, x \rangle - 1 | \leqslant \frac {\eta}{4}
\]

 对所有 \(x \in K\) 成立）；它是 \(\hat{G}\) 中 e 的一个邻域。设 \(U \subset U_{1}\) 是一个开对称邻域，它关于 Haar 测度 \(m = d\hat{x}\)（\(\hat{G}\) 的与测度 dx 对偶的 Haar 测度）可积。设 \(V \subset U\) 是 e 的一个紧对称邻域，满足 \(m(V) \geqslant \frac{1}{2}m(U)\)。用 \(\varphi_{U}\)（相应地，\(\varphi_{V}\)）表示 U（相应地，V）的特征函数。由于 \(\varphi_{U}\) 属于 \(L^{2}(G)\)，存在 \(u \in L^{2}(G)\) 使 \(\varphi_{U} = \mathcal{F}_{G}(u)\)（II，第 215 页，定理 1）。类似地，存在函数 \(v \in L^{2}(G)\) 使 \(\varphi_{V} = \mathcal{F}_{G}(v)\)。我们将证明函数 \(j = \frac{1}{m(V)}uv\) 满足所要求的性质。我们有 \(j \in L^{1}(G)\)。

\begin{enumerate}
\def\labelenumi{\alph{enumi})}
\tightlist
\item
  由 Plancherel 定理及条件 \(m(\mathrm{V}) \geqslant \frac{1}{2}m(\mathrm{U})\)，有
\end{enumerate}

\[
\| j \| _ {1} \leqslant \frac {\| u \| _ {2} \| v \| _ {2}}{m (\mathrm{V})} = \frac {\| \mathcal {F} _ {\mathrm{G}} (u) \| _ {2} \| \mathcal {F} _ {\mathrm{G}} (v) \| _ {2}}{m (\mathrm{V})} = \frac {\sqrt {m (\mathrm{U}) m (\mathrm{V})}}{m (\mathrm{V})} \leqslant \sqrt {2}.
\]

\begin{enumerate}
\def\labelenumi{\alph{enumi})}
\setcounter{enumi}{1}
\tightlist
\item
  存在 \(\widehat{G}\) 中 e 的一个邻域 W 使 WV ⊂ U（TG，II，第 31 页，命题 4）。对每个 \(\widehat{x} \in W\)，有 \(\widehat{x}V \subset U\)，且 II，第 223 页命题 14 蕴含
\end{enumerate}

\[
\begin{array}{l} \mathcal {F} _ {\mathrm{G}} (j) (\widehat {x}) = \frac {1}{m (\mathrm{V})} (\mathcal {F} _ {\mathrm{G}} (u) * \mathcal {F} _ {\mathrm{G}} (v)) (\widehat {x}) \\ \qquad = \frac {1}{m (\mathrm{V})} \int_ {\widehat {\mathrm{G}}} \varphi_ {\mathrm{U}} (\widehat {y}) \varphi_ {\mathrm{V}} (\widehat {y} ^ {- 1} \widehat {x}) d m (\widehat {y}) \\ \qquad = \frac {m (\mathrm{U} \cap \widehat {x} \mathrm{V} ^ {- 1})}{m (\mathrm{V})} = \frac {m (\widehat {x} \mathrm{V})}{m (\mathrm{V})} = 1 \end{array}
\]

因为 V 是对称的。

\begin{enumerate}
\def\labelenumi{\alph{enumi})}
\setcounter{enumi}{2}
\tightlist
\item
  若 \(x \in \mathrm{K}\)，有
\end{enumerate}

\[
\| u - u * \varepsilon_ {x} \| _ {2} ^ {2} = \int_ {\widehat {\mathrm{G}}} \big | \mathcal {F} _ {\mathrm{G}} (u) (\widehat {x}) \big (1 - \overline {{\langle x , \widehat {x} \rangle}} \big) \big | ^ {2} d m (\widehat {x}) \leqslant m (\mathrm{U}) \Big (\frac {\eta}{4} \Big) ^ {2}
\]

因为 \(\mathrm{U}\subset \mathrm{U}_1\)，同理 \(\| v - v*\varepsilon_x\| _2^2\leqslant m(\mathrm{V})\big(\frac{\eta}{4}\big)^2\)。因此

\[
\begin{array}{l} \| j - j * \varepsilon_ {x} \| _ {1} = \frac {1}{m (\mathrm{V})} \| u (v - v * \varepsilon_ {x}) + (v * \varepsilon_ {x}) (u - u * \varepsilon_ {x}) \| _ {1} \\ \leqslant \frac {\eta}{4 m (\mathrm{V})} (\| u \| _ {2} \sqrt {m (\mathrm{V})} + \| v \| _ {2} \sqrt {m (\mathrm{U})}) \\ = \frac {\eta \sqrt {m (\mathrm{U}) m (\mathrm{V})}}{2 m (\mathrm{V})} <   \eta . \end{array}
\]

命题 3.　代数 L \(^{1}\) (G) 满足 Ditkin 条件（I，第 92 页，定义 2）。

设 \(\chi\) 是 \(\mathrm{L}^1 (\mathrm{G})\) 的一个特征标。我们按 \(\chi\) 是否为零区分两种情形。若 \(\chi\) 为零，则须验证：对每个函数 \(f\in \mathrm{L}^{1}(\mathrm{G})\)，存在序列 \((f_n)_{n\geqslant 1}\)（在 \(\mathrm{L}^1 (\mathrm{G})\) 中），使 \(\mathcal{F}(f_n)\) 在 \(\widehat{\mathbf{G}}\) 的某个紧子集之外为零，且 \(f_{n}*f\) 趋于 \(f\)（在 \(\mathrm{L}^1 (\mathrm{G})\) 中）。这样的序列的存在性由上面的引理 1 及 II，第 211 页命题 8 得到。

现在设 \(\chi\) 非零，于是 \(\chi \in \mathsf{X}(\mathrm{L}^1 (\mathrm{G})) = \widehat{\mathrm{G}}\)（II，第 202 页，命题 1）。设 \(f\in \mathrm{L}^{1}(\mathrm{G})\) 满足 \(\mathcal{G}_{\mathrm{L}^1 (\mathrm{G})}(f)(\chi) = \overline{\mathcal{F}} (f)(\chi) = 0\)。剩下要证明存在序列 \((f_n)_{n\geqslant 1}\)（在

\(L^{1}(G)\) 中），使 \(f * f_{n}\) 在 \(L^{1}(G)\) 中收敛到 f，且 \(\overline{\mathcal{F}}(f_{n})\) 在 \(\chi\) 的一个邻域内为零。可设 \(\|f\|_{1}=1\)。通过 \(\widehat{G}\) 中的平移，可化归为 \(\chi=e\) 的情形。

设 \(K_{n}\) 是 G 的一个紧子集，使得

\[
\int_ {\mathrm{G} - \mathrm{K} _ {n}} | f (x) | d x \leqslant \frac {1}{n}.
\]

设 \(u_{n} \in L^{1}(G)\) 是一个 \(\geqslant 0\) 的函数，满足 \(\|u_{n}\|_{1} = 1\) 且

\[
\| f - f * u _ {n} \| _ {1} \leqslant \frac {1}{n}
\]

（参看 II，第 211 页，命题 8, (iii)）。由引理 2，存在函数 \(j_n\)（在 \(\mathrm{L}^1(\mathrm{G})\) 中），使 \(\| j_n\| _1\leqslant \sqrt{2}\)，其 Fourier 余变换在 \(e\) 的一个邻域内等于 1，并且 \(\| j_{n} - j_{n}*\varepsilon_{x}\|_{1}\leqslant n^{-1}\) 对所有 \(x\in \mathrm{K}_n\) 成立。我们置

\[
f _ {n} = u _ {n} - j _ {n} * u _ {n}.
\]

我们将证明序列 \((f_{n})_{n\geqslant1}\) 具有所要求的性质。首先，我们有

\[
\mathcal {F} (f _ {n}) = \mathcal {F} (u _ {n}) - \mathcal {F} (j _ {n}) \mathcal {F} (u _ {n}) = (1 - \mathcal {F} (j _ {n})) \mathcal {F} (u _ {n})
\]

于是 \(f_{n}\) 的 Fourier 变换在 \(\chi = e\) 的一个邻域内为零。另一方面，

\[
\| f * f _ {n} - f \| _ {1} \leqslant \| f * u _ {n} - f \| _ {1} + \| f * j _ {n} \| _ {1} \| u _ {n} \| _ {1} \leqslant \frac {1}{n} + \| f * j _ {n} \| _ {1}.
\]

现在，对几乎每个 \(y \in G\)，有

\[
(f * j _ {n}) (y) = \int_ {\mathrm{G}} f (x) j _ {n} (x ^ {- 1} y) d x = \int_ {\mathrm{G}} f (x) (j _ {n} (x ^ {- 1} y) - j _ {n} (y)) d x
\]

因为由假设有 \(\mathcal{F}(f)(e) = \int_{\mathrm{G}} f(x) dx = 0\)。由此

\[
\begin{array}{l} \| f * j _ {n} \| _ {1} \leqslant \int_ {\mathrm{G}} | f (x) | \| j _ {n} * \varepsilon_ {x} - j _ {n} \| _ {1} d x \\ \qquad = \int_ {\mathrm{K} _ {n}} | f (x) | \| j _ {n} * \varepsilon_ {x} - j _ {n} \| _ {1} d x \\ \qquad + \int_ {\mathrm{G-K} _ {n}} | f (x) | \| j _ {n} * \varepsilon_ {x} - j _ {n} \| _ {1} d x \\ \qquad \leqslant \frac {1}{n} \int_ {\mathrm{K} _ {n}} | f (x) | d x + 4 \int_ {\mathrm{G-K} _ {n}} | f (x) | d x \leqslant \frac {5}{n}. \end{array}
\]

最后，\(\|f * f_{n} - f\|_{1} \leqslant 6n^{-1}\)，因此 \(f * f_{n}\) 收敛到 \(f\)（在 \(L^{1}(G)\) 中），如愿。

应用命题 5，我们进而得到如下结果：

定理 2.　设 I 是 L \(^{1}\) (G) 的一个闭理想，使得 V(I) 的边界不含任何非空完全集。则 I 是由这样的函数 \(f \in \text{L}^{1}(\text{G})\) 所组成的集合：\(\mathcal{F}(f)\) 在 V(I) 上为零。

对 \(L^{1}(G)\) 的任意闭理想，定理 2 的结论一般不成立（参看 II，第 314 页，习题 12）。更确切地说，可以证明：若 G 非紧，则存在 \(L^{1}(G)\) 的一个非自伴的闭理想（例如见 W. RUDIN, Fourier analysis on groups, Interscience tracts in pure and applied mathematics, 定理 7.7.1.）

推论.　若 L \(^{1}\) (G) 的一个闭理想 I 只含于一个正则极大理想，则 I 本身是正则极大的。

